\documentclass[10pt,a4paper]{amsart}
\usepackage[margin=19mm]{geometry}
\usepackage{amsmath,amssymb,mathtools}
\usepackage[scr=rsfs,cal=euler]{mathalfa}
\usepackage{xfrac}
\usepackage[unicode,hidelinks,bookmarks=false]{hyperref}
\newcommand{\ie}{i.e.\ }
\newcommand{\cf}{cf.\ }
\newcommand{\resp}{respectively\ }
\newcommand{\Subsection}{\subsection}
\providecommand{\nobreakdash}{\nobreak\hbox{-}}
\setlength{\emergencystretch}{2em}
\multlinegap0pt
\usepackage{graphicx}
\usepackage{subcaption}
\usepackage[dvipsnames,svgnames,table]{xcolor}
\usepackage{tikz}
\usetikzlibrary{arrows.meta,calc,bending}

\definecolor{figgreen}{RGB}{0,112,76}
\definecolor{figred}{RGB}{184,48,38}
\definecolor{figblue}{RGB}{0,92,170}
\definecolor{figyellow}{RGB}{150,110,0}

\newcommand{\bW}{0.95}
\newcommand{\bStep}{1.82}
\newcommand{\rowGap}{1.60}
\newcommand{\padA}{0.40}
\newcommand{\inSep}{0.45}
\newcommand{\grGap}{0.85}
\newcommand{\padB}{0.33}

\tikzset{
  fgdot/.style   ={circle,draw=none,fill,inner sep=0pt,minimum size=4.6pt},
  gdot/.style    ={fgdot,fill=figgreen},
  rdot/.style    ={fgdot,fill=figred},
  kdot/.style    ={fgdot,fill=black},
  fgbox/.style   ={draw=black,line width=0.8pt},
  fgline/.style  ={line width=1.15pt,line cap=round},
  fgarr/.style   ={fgline,-{Stealth[length=4.6pt,width=4.4pt,bend]},shorten >=1.6pt},
  gedge/.style   ={fgarr,figgreen},
  redge/.style   ={fgarr,figred,dash pattern=on 4pt off 2pt},
  kedge/.style   ={fgarr,black,line width=1.0pt},
  bedge/.style   ={fgarr,figblue,dash pattern=on 1.4pt off 1.2pt},
  yedge/.style   ={fgarr,figyellow,dash pattern=on 5pt off 1.5pt on 1pt off 1.5pt},
  glab/.style    ={inner sep=1pt,fill=white},
  dts/.style     ={inner sep=1pt,font=\bfseries},
}
\newcommand{\dt}{$\cdots$}

\pgfdeclarelayer{fglabels}
\pgfsetlayers{main,fglabels}
\newcommand{\toplabel}[3]{\begin{pgfonlayer}{fglabels}\node[glab] at (#1,#2) {#3};\end{pgfonlayer}}

\newcommand{\grbox}[4]{%
  \pgfmathsetmacro{\bhh}{(\rowGap+2*\padA)/2}%
  \pgfmathsetmacro{\bww}{\bW/2}%
  \pgfmathsetmacro{\ytop}{#3+\rowGap/2}%
  \pgfmathsetmacro{\ybot}{#3-\rowGap/2}%
  \draw[fgbox] (#2-\bww,#3-\bhh) rectangle (#2+\bww,#3+\bhh);
  \toplabel{#2}{#3}{#4}%
  \node[gdot] (#1i1) at (#2-\bww,\ytop) {};
  \node[gdot] (#1o1) at (#2+\bww,\ytop) {};
  \node[rdot] (#1i2) at (#2-\bww,\ybot) {};
  \node[rdot] (#1o2) at (#2+\bww,\ybot) {};
}
\newcommand{\grboxx}[3]{\grbox{#1}{#2}{#3}{}}

\newcommand{\nbox}[6]{%
  \pgfmathsetmacro{\nk}{#4+#5}%
  \pgfmathsetmacro{\htot}{(#4-1)*\inSep+(#5-1)*\inSep+\grGap}%
  \pgfmathsetmacro{\bhh}{(\htot+2*\padB)/2}%
  \pgfmathsetmacro{\bww}{\bW/2}%
  \pgfmathsetmacro{\ylab}{#3+\htot/2-(#4-1)*\inSep-\grGap/2}%
  \draw[fgbox] (#2-\bww,#3-\bhh) rectangle (#2+\bww,#3+\bhh);
  \toplabel{#2}{\ylab}{#6}%
  \foreach \i in {1,...,#4}{%
    \pgfmathsetmacro{\yy}{#3+\htot/2-(\i-1)*\inSep}%
    \node[kdot] (#1i\i) at (#2-\bww,\yy) {};
    \node[kdot] (#1o\i) at (#2+\bww,\yy) {};
  }%
  \foreach \i in {1,...,#5}{%
    \pgfmathsetmacro{\yy}{#3-\htot/2+(#5-\i)*\inSep}%
    \pgfmathtruncatemacro{\j}{#4+\i}%
    \node[kdot] (#1i\j) at (#2-\bww,\yy) {};
    \node[kdot] (#1o\j) at (#2+\bww,\yy) {};
  }%
}
\newcommand{\nboxx}[5]{\nbox{#1}{#2}{#3}{#4}{#5}{}}

\newcommand{\rowy}[4]{%
  \pgfmathparse{(#4<=#2) ? (#1+((#2-1)*\inSep+(#3-1)*\inSep+\grGap)/2-(#4-1)*\inSep)
                         : (#1-((#2-1)*\inSep+(#3-1)*\inSep+\grGap)/2+(#2+#3-#4)*\inSep)}\pgfmathresult}

\newcommand{\otimescircle}[3][0.30]{%
  \draw[line width=1.0pt] (#2,#3) circle (#1);
  \draw[line width=1.0pt] (#2-0.707*#1,#3-0.707*#1) -- (#2+0.707*#1,#3+0.707*#1);
  \draw[line width=1.0pt] (#2-0.707*#1,#3+0.707*#1) -- (#2+0.707*#1,#3-0.707*#1);
}

\newcommand{\uparc}[4][gedge]{%
  \draw[#1] (#2) .. controls ($(#2)+(0.30,#4)$) and ($(#3)+(-0.30,#4)$) .. (#3);}
\newcommand{\dnarc}[4][redge]{%
  \draw[#1] (#2) .. controls ($(#2)+(0.30,-#4)$) and ($(#3)+(-0.30,-#4)$) .. (#3);}
\newcommand{\uparcs}[5][gedge]{%
  \draw[#1] (#2) .. controls ($(#2)+(#5,#4)$) and ($(#3)+(-#5,#4)$) .. (#3);}
\newcommand{\dnarcs}[5][redge]{%
  \draw[#1] (#2) .. controls ($(#2)+(#5,-#4)$) and ($(#3)+(-#5,-#4)$) .. (#3);}

\newcommand{\rgbox}[4]{%
  \pgfmathsetmacro{\bhh}{(\rowGap+2*\padA)/2}%
  \pgfmathsetmacro{\bww}{\bW/2}%
  \pgfmathsetmacro{\ytop}{#3+\rowGap/2}%
  \pgfmathsetmacro{\ybot}{#3-\rowGap/2}%
  \draw[fgbox] (#2-\bww,#3-\bhh) rectangle (#2+\bww,#3+\bhh);
  \toplabel{#2}{#3}{#4}%
  \node[rdot] (#1i1) at (#2-\bww,\ytop) {};
  \node[rdot] (#1o1) at (#2+\bww,\ytop) {};
  \node[gdot] (#1i2) at (#2-\bww,\ybot) {};
  \node[gdot] (#1o2) at (#2+\bww,\ybot) {};
}

\newsavebox{\ffbox}
\newlength{\ffW}\newlength{\ffH}\newlength{\ffT}
\newcommand{\fitfig}[3]{%
  \setlength{\ffW}{#1}\setlength{\ffH}{#2}%
  \sbox{\ffbox}{#3}%
  \setlength{\ffT}{\dimexpr\ht\ffbox+\dp\ffbox\relax}%
  \pgfmathsetmacro{\ffs}{min(\ffW/\wd\ffbox,\ffH/\ffT)}%
  \begin{minipage}[b][\ffH][c]{\ffW}%
    \centering\scalebox{\ffs}{\usebox{\ffbox}}%
  \end{minipage}%
}

\emergencystretch=2em
%\usepackage{showkeys}

\allowdisplaybreaks[4]

\numberwithin{equation}{section}

\theoremstyle{plain}
\newtheorem{theorem}{Theorem}[section]
\newtheorem{lemma}[theorem]{Lemma}
\newtheorem{corollary}[theorem]{Corollary}
\newtheorem{proposition}[theorem]{Proposition}
\newtheorem{conjecture}[theorem]{Conjecture}
\newtheorem{assumption}[theorem]{Assumption}
\newtheorem{example}[theorem]{Example}
\newtheorem{definition}[theorem]{Definition}
\newtheorem{remark}[theorem]{Remark}
\newtheorem*{theorem*}{Theorem}

\def\bE{\mathbb E}
\def\bP{\mathbb P}
\def\bN{\mathbb N}
\def\bR{\mathbb R}
\def\bC{\mathbb C}
\def\bM{\mathbb M}

\def\cA{\mathcal A}
\def\cB{\mathcal B}
\def\cG{\mathcal G}
\def\cP{\mathcal P}

\def\red{\color{red}}
\def\blue{\color{blue}}

\def\iunit{\mathrm {i}}
\def\id{\mathrm {Id}}
\def\Tr{\mathrm {Tr}}
\def\tr{\mathrm {tr}}
\def\diag{\mathrm {diag}}
\def\Var{\mathrm {Var}}
\def\Cov{\mathrm {Cov}}
\def\free{\mathrm {free}}
\def\sa{\mathrm {sa}}
\def\Wg{\mathrm {Wg}}
\title[Operator norm bounds for multi-leg matrix tensors]{Operator Norm Bounds for Multi-leg Matrix Tensors and Applications to Random Matrix Theory: the two-leg scalar bound (Theorem 3.1)}
\author[Benoît Collins and Wangjun Yuan]{Benoît Collins and Wangjun Yuan}
\date{Source: arXiv:2603.27659v2 (CC BY 4.0)}
\begin{document}
\begin{abstract}
We study the extremal values of multi-leg traces of matrix tensors under operator norm constraints. A graphical representation gives upper and lower bounds expressed as powers of the common tensor-factor dimension. The lower bounds are attained by matrices that permute tensor factors and are exact within this family. We prove the bounds first for two tensor factors and then for an arbitrary number. Leaving some indices uncontracted produces matrices for which we obtain both moment and operator norm bounds; a single choice of coefficient matrices attains the lower bounds for every positive integer moment. We also obtain exact scalar and operator norm maxima in special cases with sufficiently many tensor factors sharing the same cyclic contraction. As an application, the operator norm bounds yield a uniform comparison between Ginibre products and their free circular counterparts with growing matrix coefficients.
\end{abstract}

\maketitle
\vspace{1em}
\noindent\textit{Source and licence.} Operator Norm Bounds for Multi-leg Matrix Tensors and Applications to Random Matrix Theory. Version arXiv:2603.27659v2 (CC BY 4.0), distributed under the Creative Commons Attribution 4.0 International licence, http://creativecommons.org/licenses/by/4.0/, as recorded in the arXiv OAI licence metadata for this exact version and shown on the abstract page https://arxiv.org/abs/2603.27659v2. This item is an editorial re-presentation of the paper's Theorem 3.1 with the paper's own proof of it, together with the source-local chain that proof uses. A full rights notice is delivered at licenses/arxiv-2603.27659-CC-BY-4.0.txt.
\noindent\textit{Editorial scope.} Counted unit: the paper's two-leg scalar operator-norm bound, printed as Theorem 3.1 (source0.tex lines 706--714), delivered together with the paper's complete original proof of it (source0.tex lines 1474--1481) as the single primary proof body. No mathematics is written, repaired or re-derived: statement and proof are the paper's own text line for line, in the paper's own notation ($\Tr_{\sigma_1}\otimes\Tr_{\sigma_2}$, $\cP([m])$, $M(\sigma_1,\sigma_2)$, $C(\sigma_1,\sigma_2)$, $G_{\sigma_1,\sigma_2}$, $F_\prec$, $U$). Required context retained uncounted, because the counted proof is the paper's final assembly over it: the abstract (source0.tex lines 189--191, reproduced in the wrapper front matter); the graphical notation of the paper's Section 2 with its coordinate trace, its partial trace, its graph for two legs and its Definition 2.1 of simple and full partial graphs (source0.tex lines 262--263, 266--275, 280--297 and 416--440); the paper's Section 3 statement block with Theorem 3.1, its two corollaries, the optimal-ordering proposition with its own proof and the subsequent remark (lines 698--760); the whole of the paper's Section 4 with Theorem 4.1, the contraction identity and Lemmas 4.2--4.4 with their complete proofs, including the long induction of Lemma 4.4 (lines 912--950, 1012--1372 and 1374--1404); and the whole of the paper's Sections 5 and 6 with Theorem 5.1 and its proof and the paper's own proofs of Theorem 3.1 and its contraction corollary (lines 1406--1481). Printed numbers are the paper's own: the wrapper restores the paper's section counter before the first retained section, and keeps the paper's own section-relative theorem counter and section-relative equation counter, so the retained sections print as Sections 2 to 6, the counted theorem as Theorem 3.1, the retained definitions, lemmas, propositions and corollaries with the paper's numbers, and the numbered displays as the paper's (2.1)--(6.1). Omitted from this package and not redistributed: the paper's preamble and front matter as such (lines 1--61 and 177--198, of which only the title, the author list and the abstract are reproduced in the wrapper), the introduction and literature review (lines 200--260), the parts of the paper's Section 2 outside the retained ranges, the multiple-leg results of the paper's Section 7 and its applications section, and the paper's own bibliography machinery with its bib data file. No citation command occurs in any retained span, so the delivered document carries no bibliography, and no reference command of the delivered text was rewritten. Scope deviation, disclosed: the counted proof is the paper's short final assembly (the lower bound is taken from the retained Theorem 5.1, the upper bound from the retained Theorem 4.1); the difficulty of the unit is the retained chain, which is delivered in full precisely so that the counted assembly can be checked: the graphical calculus, the Cauchy--Schwarz pairing of a full simple partial graph with its complement, the Hilbert--Schmidt norm computation for those graphs, and the permutation-tensor lower bound. Printed slips kept literal, among them the paper's own wording and its own choice of symbols. Layout and class: the paper's own class is the standard article class at 12pt, which is not a publisher class, and it is not shipped with this package; the delivered wrapper is the lane's pinned amsart editorial wrapper at 10pt on A4, to which this package adds the paper's own TikZ drawing environment with its colours, styles and drawing macros, its own ten theorem-like declarations, its section-relative equation numbering and its own mathematical shorthand macros, all copied verbatim from the paper's preamble. No publisher class file, no publisher style file and no upstream script is redistributed. Assets: the paper's computations are drawn inline with its own TikZ code; six figure environments and their TikZ drawings from the retained spans are reproduced as the paper's own source code, no external image file exists in the package, no retained span carries a graphics-inclusion command, and the manifest and the delivery evidence both carry an empty asset-dependency list. A full rights notice is delivered at licenses/arxiv-2603.27659-CC-BY-4.0.txt.
\setcounter{section}{1}
\section{Graphical notation} \label{sec:graph}


We use the standard basis of $(\bC^N)^{\otimes k}$. Write
$(A_i)_{t_i^1\cdots t_i^k;s_i^1\cdots s_i^k}$ for the tensor matrix entries of $A_i$, where all indices range over $[N]$. The row indices $t$ label in-vertices and the column indices $s$ label out-vertices. For permutations, the index contraction is
\begin{equation}\label{eq-coordinate-trace}
(\Tr_{\sigma_1}\otimes\cdots\otimes\Tr_{\sigma_k})(A_1,\ldots,A_m)
=\sum_{\mathbf t,\mathbf s}
\prod_{i=1}^m (A_i)_{t_i^1\cdots t_i^k;s_i^1\cdots s_i^k}
\prod_{j=1}^k\prod_{i=1}^m\delta_{s_i^j,t_{\sigma_j(i)}^j},
\end{equation}
where $\sum_{\mathbf t,\mathbf s}$ sums over $t_i^l$ and $s_i^l$ for all $l \in [k]$ and $i \in [m]$.
Thus each colored edge identifies its two endpoint indices.

\subsection{Graph for 2 legs} \label{sec:graph-2 leg}

We start with the case $k=2$ as a warm-up. In this manuscript, we take the stance that a good understanding of the case with two legs will facilitate the statement and the proof of the general case.
Our first goal is to evaluate and bound this contraction: for permutations $\sigma_1, \sigma_2 \in \cP([m])$, we compute the following partial trace:
\begin{align} \label{tr-sigma-tau}
    \left( \Tr_{\sigma_1} \otimes \Tr_{\sigma_2} \right) \left( A_1, \ldots, A_m \right)
\end{align}
and bound its maximum among all unitary matrices $A_1, \ldots, A_m \in M_N (\bC) \otimes M_N(\bC)$.

We use a graph with the notation $G_{\sigma_1,\sigma_2} (A_1,\ldots,A_m)$ for the representation of the partial trace \eqref{tr-sigma-tau}. We use $m$ rectangles for the matrices $A_1,\ldots,A_m$. We connect the rectangles with directed edges according to the permutations $\sigma_1, \sigma_2$.

For the permutation $\sigma_1$, for all $i$, we use a green directed edge to connect the rectangles $A_i$ and $A_{\sigma_1(i)}$ with the orientation from the rectangle $A_i$ to the rectangle $A_{\sigma_1(i)}$. The vertices of the green edge on the rectangle $A_i$ and the rectangle $A_{\sigma_1(i)}$ are called the green \textit{out-vertex} and the green \textit{in-vertex}, respectively. We also call the green edge an \textit{out-edge} of $A_i$ and an \textit{in-edge} of $A_{\sigma_1(i)}$.
We use red directed edges to connect rectangles for any permutation $\sigma_2$, following the same rule.
That is, the red directed edges are from $A_i$ to $A_{\sigma_2(i)}$.
Then, every rectangle has one green in-edge, one green out-edge, one red in-edge, and one red out-edge. Thus, there are four vertices on every rectangle: the green in-vertex, the green out-vertex, the red in-vertex, and the red out-vertex.
Moreover, if we shrink every rectangle to a point, then the connected components of the red and green directed edges correspond to the cycles of the permutations $\sigma_2$ and $\sigma_1$, respectively.
It is obvious that the graph defined above has $m$ red edges and $m$ green edges.
For simplicity, unless otherwise specified, we always draw the out-vertices on the right of the rectangles, and the in-vertices on the left of the rectangles. Unless specified otherwise, the quantities we consider in the sequel do not depend on the position of each rectangle. Without loss of generality, we always list the rectangles $A_1,\ldots,A_m$ from left to right.

\begin{definition}\label{def-cycle-cut}
Let $G_{\sigma_1,\ldots,\sigma_k}(A_1,\ldots,A_m)$ be a graph of permutations $\sigma_1,\ldots,\sigma_k$.
For a blue-edge configuration $\boldsymbol\pi=(\pi_1,\ldots,\pi_m)$, where $\pi_i(j)$ is the color of the out-vertex paired with the in-vertex of color $j$ in rectangle $A_i$, let $M(\sigma_1,\ldots,\sigma_k;\boldsymbol\pi)$ be the number of directed cycles.
Define
\[
M(\sigma_1,\ldots,\sigma_k)=\max_{\boldsymbol\pi\in\cP([k])^m}M(\sigma_1,\ldots,\sigma_k;\boldsymbol\pi).
\]
A partial graph is \emph{simple} if it contains no directed cycle for any blue-edge configuration, and \emph{full} if it retains all rectangles.

Define
\[
C(\sigma_1,\ldots,\sigma_k)=\min_{G'} R_{G'},
\]
where $\min_{G'}$ is taken over all full simple partial graphs $G'$ of $G_{\sigma_1,\ldots,\sigma_k}$, and $R_{G'}$ is the number of colored edges removed from $G_{\sigma_1,\ldots,\sigma_k}$ to obtain $G'$. We also write $C(G)$ for $C(\sigma_1,\ldots,\sigma_k)$ when $G=G_{\sigma_1,\ldots,\sigma_k}$.
\end{definition}
The minimum exists: removing all colored edges leaves a full simple partial graph. To obtain a full simple partial graph, one must cut each cycle of every fixed blue-edge configuration. Thus, $M(\sigma_1,\ldots,\sigma_k)\le C(\sigma_1,\ldots,\sigma_k)$.

The permutation-tensor associated with $\pi\in\cP([k])$ is
\begin{equation}\label{eq-permutation-tensor}
(E_\pi)_{t_1\cdots t_k;s_1\cdots s_k}
=\prod_{j=1}^k\delta_{t_j,s_{\pi(j)}}.
\end{equation}
It is unitary and realizes precisely the blue edges $j\to\pi(j)$. Substituting these entries into \eqref{eq-coordinate-trace} identifies all indices on each directed cycle, with one independent index per cycle. Consequently,
\begin{equation}\label{eq-coordinate-cycles}
(\Tr_{\sigma_1}\otimes\cdots\otimes\Tr_{\sigma_k})(E_{\pi_1},\ldots,E_{\pi_m})
=N^{M(\sigma_1,\ldots,\sigma_k;\boldsymbol\pi)}.
\end{equation}

\section{Main results} \label{sec:results}

In this paper, we obtain a result irrespective of the number of legs in the tensor model. However, for the sake of clarity, we choose to ``warm up'' with the two leg case, whose proof we give in detail, before tackling the general $k$-legs model.

\subsection{Partial trace with 2 legs}

Our first result gives upper and lower bounds for the maximum of the partial trace under operator norm constraints.


\begin{theorem} \label{Thm-main}
Let $m,N\in\bN$ and $\sigma_1,\sigma_2\in\cP([m])$. Let $M(\sigma_1,\sigma_2)$ and $C(\sigma_1,\sigma_2)$ be the quantities given in Definition~\ref{def-cycle-cut}. Then
\begin{align*}
    N^{M(\sigma_1,\sigma_2)} \le
    \max_{A_1,\ldots,A_m} \left| \left( \Tr_{\sigma_1} \otimes \Tr_{\sigma_2} \right) \left( A_1, \ldots, A_m \right) \right|
    \le N^{C(\sigma_1,\sigma_2)},
\end{align*}
where the maximum is taken among all unitary matrices $A_1, \ldots, A_m \in M_N(\bC) \otimes M_N(\bC)$.
\end{theorem}


The same bounds hold for contractions, by convexity and multilinearity. The reduction is proved in Section~\ref{sec:2legs}.

\begin{corollary} \label{Coro-main-2 leg}
For $m,N \in \bN$, for any permutations $\sigma_1,\sigma_2 \in \cP([m])$, we have
\begin{align*}
    N^{M(\sigma_1,\sigma_2)} \le
    \max_{\|A_1\|,\ldots,\|A_m\| \le 1} \left| \left( \Tr_{\sigma_1} \otimes \Tr_{\sigma_2} \right) \left( A_1, \ldots, A_m \right) \right|
    \le N^{C(\sigma_1,\sigma_2)}.
\end{align*}
\end{corollary}

We next express $C(\sigma_1,\sigma_2)$ in terms of orders of the rectangles. For a partial permutation $\sigma$, write
\[
R(\sigma)=\#\{i\in D(\sigma):\sigma(i)\le i\}.
\]
In the natural order $A_1,\ldots,A_m$, a colored edge from $A_i$ to $A_{\sigma(i)}$ is called \emph{backward} if $\sigma(i)\le i$; loops are included. More generally, for a linear order $\prec$ of the rectangles, let $F_\prec$ be the set of backward colored edges in that order, namely edges whose target is at or to the left of their source when the rectangles are drawn in that order.

\begin{proposition}[Optimal ordering]\label{Prop-cut-order}
For any graph $G$ of partial permutations $\sigma_1,\ldots,\sigma_k$,
\begin{equation}
C(\sigma_1,\ldots,\sigma_k)=\min_\prec |F_\prec|.
\end{equation}
\end{proposition}
\begin{proof}
For a graph of partial permutations, collapse each rectangle to a vertex, keeping every retained colored directed edge. The resulting directed graph has no directed cycle if and only if the graph is simple.
Indeed, a cycle with blue directed edges gives a directed closed walk in the collapsed graph, which contains a directed cycle. Conversely, any directed cycle (including loops) in the collapsed graph can be realized by pairing its in-vertex and out-vertex inside each rectangle on that cycle and completing all other blue pairings arbitrarily.

Deleting all directed edges in $F_\prec$ leaves only edges going strictly forward, so $C(\sigma_1,\ldots,\sigma_k)\le |F_\prec|$ for every order. Conversely, choose a set $F$ of $C(\sigma_1,\ldots,\sigma_k)$ edges whose removal leaves a full simple partial graph. Its collapsed graph does not have any cycles and therefore has an order $\prec$ of rectangles in which every retained edge goes strictly forward. For all backward edges in this order $\prec$ of rectangles, their corresponding directed edges in the original graph belong to $F$, giving $|F_\prec|\le C(\sigma_1,\ldots,\sigma_k)$. This proves the identity.
\end{proof}

In the natural order for two legs, $|F_\prec|=R(\sigma_1)+R(\sigma_2)$. The upper bound in Corollary~\ref{Coro-main-2 leg} therefore gives the following estimate.

\begin{corollary} \label{Coro-2 leg-backward}
For $m,N \in \bN$, for any permutations $\sigma_1,\sigma_2 \in \cP([m])$, we have
\begin{align*}
    \max_{\|A_1\|,\ldots,\|A_m\| \le 1} \left| \left( \Tr_{\sigma_1} \otimes \Tr_{\sigma_2} \right) \left( A_1, \ldots, A_m \right) \right|
    \le N^{R(\sigma_1) + R(\sigma_2)}.
\end{align*}
\end{corollary}

\begin{remark}
The backward-edge count depends on the order of the rectangles, whereas the maximum of the index contraction is invariant under simultaneous relabelling. Thus, for every $\theta\in\cP([m])$,
\[
\max_{\|A_i\|\le1}|(\Tr_{\sigma_1}\otimes\Tr_{\sigma_2})(A_1,\ldots,A_m)|
\le N^{R(\theta\sigma_1\theta^{-1})+R(\theta\sigma_2\theta^{-1})}.
\]
For example, if $\sigma_1=(164253)$ and $\sigma_2=(123456)$, then $R(\sigma_1)=4$ and $R(\sigma_2)=1$, giving $N^5$. Relabelling by $\theta=(15)(26)$ gives $\theta\sigma_1\theta^{-1}=(524613)$ and $\theta\sigma_2\theta^{-1}=(563412)$, each with two backward edges, and improves the bound to $N^4$.

By Proposition~\ref{Prop-cut-order}, optimizing over all relabellings gives exactly $N^{C(\sigma_1,\sigma_2)}$.
\end{remark}

\section{Upper bound in the 2 legs case} \label{sec:upper}

In this section, we establish the upper bound for the partial trace \eqref{tr-sigma-tau}.

\subsection{Upper bound for partial trace}

For $\sigma_1,\sigma_2\in\cP([m])$, consider the graph $G_{\sigma_1,\sigma_2}(A_1,\ldots,A_m)$ defined in Section~\ref{sec:graph-2 leg}.

Recall the notions of simple and full partial graphs from Definition~\ref{def-cycle-cut}. If $G'_{\sigma_1,\sigma_2}$ is full, its complement $G''_{\sigma_1,\sigma_2}$ consists only of the removed colored edges, whose number is $R_{G'_{\sigma_1,\sigma_2}}$.

The following result is the upper bound for the partial trace \eqref{tr-sigma-tau}, and is the main result in this section.

\begin{theorem} \label{Thm-upper}
For any $m\in\bN$ and permutations $\sigma_1,\sigma_2 \in \cP([m])$, for any $N \in \bN$, we have
\begin{align*}
    \max_{A_1,\ldots,A_m} \left| \left( \Tr_{\sigma_1} \otimes \Tr_{\sigma_2} \right) \left( A_1, \ldots, A_m \right) \right|
    \le N^{C(\sigma_1,\sigma_2)}.
\end{align*}
Here, $\max_{A_1,\ldots,A_m}$ is taken over all unitary matrices $A_1,\ldots,A_m \in M_N (\bC) \otimes M_N(\bC)$.
\end{theorem}

\subsection{Proof of Theorem \ref{Thm-upper}}

We evaluate a partial graph as a tensor using the coordinate convention of Section~\ref{sec:graph}. Each retained rectangle contributes its matrix entry, and each retained colored edge contributes a Kronecker delta. An index of the removed directed edge remains free: on a retained rectangle it labels an open in-vertex or out-vertex, while on the complementary tensor of removed edges it labels the corresponding free vertex. Sum over all indices whose corresponding vertices are not open, and leave these free indices unsummed. With any fixed ordering of the free indices, define
\[
\langle X,Z\rangle=\sum_{\boldsymbol\alpha}X_{\boldsymbol\alpha}\overline{Z_{\boldsymbol\alpha}},
\qquad \|X\|_2^2=\sum_{\boldsymbol\alpha}|X_{\boldsymbol\alpha}|^2.
\]
This is the Hilbert--Schmidt pairing after grouping endpoints into row and column indices.

\begin{theorem*}[Contraction identity]
For permutations $\sigma_1,\sigma_2\in\cP([m])$ and the corresponding graph $G_{\sigma_1,\sigma_2}(A_1,\ldots,A_m)$, let $G'_{\sigma_1,\sigma_2}$ be a full partial graph and $G''_{\sigma_1,\sigma_2}$ be the complement of $G'_{\sigma_1,\sigma_2}$, with its free endpoint indices ordered to match the open vertex indices of $G'_{\sigma_1,\sigma_2}$. Then
\[
G_{\sigma_1,\sigma_2} =\langle G'_{\sigma_1,\sigma_2}, G''_{\sigma_1,\sigma_2} \rangle.
\]
\end{theorem*}
\begin{proof}
The tensor $G''_{\sigma_1,\sigma_2}$ is a product of Kronecker deltas and is real. Summing the product of the entries of $G'_{\sigma_1,\sigma_2}$ and $G''_{\sigma_1,\sigma_2}$ over their common endpoint indices restores exactly the removed identifications in \eqref{eq-coordinate-trace}. Complex conjugation on $G''_{\sigma_1,\sigma_2}$ has no effect, which gives the stated pairing.
\end{proof}

Next, we introduce the properties of simple partial graphs used in the norm computation. In Lemmas~\ref{Lem-simple}--\ref{Lem-inner product}, we consider partial graphs whose retained directed edges have both endpoints on retained rectangles. Thus their free indices are the open in-vertices and out-vertices of those rectangles. This class includes every full partial graph used in the upper bound and is preserved when a rectangle and all its incident edges are removed. Edge-only complementary tensors are evaluated separately.

\begin{lemma} \label{Lem-simple-subgraph}
For any $m \in \bN$, for any permutations $\sigma_1,\sigma_2 \in \cP([m])$, we consider the graph $G_{\sigma_1,\sigma_2}$. Let $H$ be a simple partial graph of the graph $G_{\sigma_1,\sigma_2}$. For any partial graph $H'$ of $H$, it is a simple partial graph of the graph $G_{\sigma_1,\sigma_2}$.
\end{lemma}

This lemma follows directly from the definition of a simple partial graph.

\begin{lemma} \label{Lem-simple}
For any $m \in \bN$, for any permutations $\sigma_1,\sigma_2 \in \cP([m])$, we consider the graph $G_{\sigma_1,\sigma_2}$. Let $H$ be a simple partial graph of the graph $G_{\sigma_1,\sigma_2}$, such that the endpoints of all directed edges are vertices of rectangles. Then in $H$, there exists a rectangle $A_i$, such that either $A_i$ does not have out-edges or $A_i$ does not have in-edges.
\end{lemma}

\begin{proof}
We only consider the graph $H$ that contains at least one rectangle.
We prove the lemma by contradiction. Assume that all rectangles in $H$ have in-edges and out-edges. We will connect the blue directed edges in the rectangles to construct a directed cycle in the following way.

If there is a loop in $H$, then we connect the in-vertex and the out-vertex of this loop with a blue directed edge to form a directed cycle. In the following, we only consider the case where there is no loop in $H$.

We can start with a rectangle $A_{i_1}$ for $i_1 \in [m]$. By assumption, the rectangle $A_{i_1}$ has an out-edge. We denote by $A_{i_2}$ the successor of $A_{i_1}$ with respect to this out-edge, where $i_2 \in [m]$ and $i_2 \not= i_1$.
Since the rectangle $A_{i_2}$ also has an out-edge, we denote by $A_{i_3}$ the successor of $A_{i_2}$ with respect to this out-edge, where $i_3 \in [m]$ and $i_3 \not= i_2$.
We continue this procedure to obtain a sequence of indices $i_1 \not= i_2 \not= \ldots$, such that there is a directed edge from rectangle $A_{i_r}$ to rectangle $A_{i_{r+1}}$ for $r \in \bN$.
As the number of rectangles in the graph $H$ is finite, the sequence of indices $i_1 \not= i_2 \not= \ldots$ contains only finitely many different numbers. Hence, there exist integers $t < s$, such that $i_t=i_s$, and $i_t, i_{t+1}, \ldots, i_{s-1}$ are distinct.
For every $t<r<s$, we use a blue directed edge in rectangle $A_{i_r}$ to connect the in-vertex of the in-edge from $A_{i_{r-1}}$ to $A_{i_r}$ with the out-vertex of the out-edge from $A_{i_r}$ to $A_{i_{r+1}}$.
We provide an example of the blue directed edges in Figure \ref{figure-Lemma-cycle-blue}.
For the repeated rectangle $A_{i_t} = A_{i_s}$, we connect the blue directed edge from the in-vertex on $A_{i_s}$ to the out-vertex on $A_{i_t}$ to close the directed cycle $A_{i_t}\to A_{i_{t+1}}\to\cdots\to A_{i_s}=A_{i_t}$ and
complete each local pairing arbitrarily.

\begin{figure}[ht]
    \centering
    \fitfig{229.01bp}{80.00bp}{%
    \begin{tikzpicture}
      \useasboundingbox (-1.982,-1.433) rectangle (6.411,1.285);
      \grbox{A}{0}{0}{$A_{i_{r-1}}$}
      \grbox{B}{2.42}{0}{$A_{i_{r}}$}
      \grbox{C}{4.73}{0}{$A_{i_{r+1}}$}
      \draw[bedge] (Ai1) -- (Ao2);  \draw[bedge] (Ai2) -- (Ao1);
      \draw[bedge] (Bi1) -- (Bo2);  \draw[bedge] (Bi2) -- (Bo1);
      \draw[bedge] (Ci1) -- (Co1);  \draw[bedge] (Ci2) -- (Co2);
      \draw[gedge] (-1.94,0.80) .. controls (-1.44,1.38) and (-0.98,1.38) .. (Ai1);
      \uparcs[gedge]{Bo1}{Ci1}{0.58}{0.50}
      \draw[gedge] (Co1) .. controls (5.71,1.38) and (5.93,1.38) .. (6.43,0.80);
      \dnarcs[redge]{Ao2}{Bi2}{0.75}{0.40}
      \foreach \x in {-1.215,1.094,3.739,5.842}{
        \node[dts] at (\x,0.80) {\dt};
        \node[dts] at (\x,-0.80) {\dt};
      }
    \end{tikzpicture}%
  }
    \caption{Blue directed edges for cycles}
    \label{figure-Lemma-cycle-blue}
\end{figure}

Since $H$ is simple, the directed cycle constructed by these blue edges is a contradiction. This concludes the proof.
\end{proof}

In order to establish the upper bound, we need the following lemmas to compute the norm of the simple partial graph.

\begin{lemma} \label{Lem-inner product}
For any $m \in \bN$, for any permutations $\sigma_1,\sigma_2 \in \cP([m])$, we consider the graph $G_{\sigma_1,\sigma_2}$. Assume that all matrices labeling rectangles are unitary. Let $H$ be a simple partial graph of the graph $G_{\sigma_1,\sigma_2}$ such that the endpoints of all directed edges are vertices of rectangles. Let $n$ be the number of rectangles in the graph $H$, $r$ be the number of green directed edges in $H$, and $s$ be the number of red directed edges in $H$. Then we have
\begin{align*}
    \|H\|_2^2 = N^{2n-r-s}.
\end{align*}
\end{lemma}

\begin{proof}
We prove this by induction on $n$, allowing $r,s\ge0$.

We start with the case $n=1$. Note that a simple partial graph with only one rectangle must be a rectangle without edges. Thus, the graph of $H$ consists of one rectangle with no edges. In this case, we have $n=1, r=s=0$. See Figure \ref{figure-k=1-1}. The quantity $\|H\|_2^2$ is represented in Figure \ref{figure-k=1-2}. Noting that $A_1^*A_1 = I_N \otimes I_N$, we have $\|H\|_2^2 = \Tr (A_1^*A_1) =  N^2$.
\begin{figure}[ht]
    \centering
    \begin{subfigure}{0.48\textwidth}
        \centering
        \fitfig{40.80bp}{84.60bp}{%
    \begin{tikzpicture}
      \useasboundingbox (-0.600,-1.237) rectangle (0.578,1.235);
      \grbox{A}{0}{0}{$A_1$}
    \end{tikzpicture}%
  }
        \caption{Graph of $H$}
        \label{figure-k=1-1}
    \end{subfigure}
    \begin{subfigure}{0.48\textwidth}
        \centering
        \fitfig{65.00bp}{157.51bp}{%
    \begin{tikzpicture}
      \useasboundingbox (-1.108,-2.756) rectangle (1.102,2.756);
      \rgbox{U}{0}{1.52}{$A_1^{*}$}
      \grbox{L}{0}{-1.52}{$A_1$}
      \draw[gedge] (Ui2) .. controls ($(Ui2)+(-0.29,-0.50)$) and ($(Li1)+(-0.29,0.50)$) .. (Li1);
      \draw[gedge] (Lo1) .. controls ($(Lo1)+(0.29,0.50)$)  and ($(Uo2)+(0.29,-0.50)$) .. (Uo2);
      \draw[redge] (Ui1) .. controls ($(Ui1)+(-0.76,-1.40)$) and ($(Li2)+(-0.76,1.40)$) .. (Li2);
      \draw[redge] (Lo2) .. controls ($(Lo2)+(0.76,1.40)$)  and ($(Uo1)+(0.76,-1.40)$) .. (Uo1);
    \end{tikzpicture}%
  }
        \caption{Graph interpretation of $\|H\|_2^2$}
        \label{figure-k=1-2}
    \end{subfigure}
    \caption{Case $n=1$}

\end{figure}

Next, we assume that the conclusion holds for the case $n-1$. That is, for any $0\le r,s \le n-1$, for any simple graph $H'$ that has $n-1$ rectangles with $r$ green directed edges and $s$ red directed edges, the conclusion holds.
Now we consider the graph $H$ with $n$ rectangles $A_1,\ldots,A_n$. Noting that $H$ is a simple partial graph of $G_{\sigma_1,\sigma_2}$, by Lemma \ref{Lem-simple}, there exists a rectangle which either does not have out-edges or does not have in-edges. Without loss of generality, we assume that the rectangle $A_1$ does not have any in-edges. We consider the number of out-edges of $A_1$. We have the following four cases.

\bigskip
\noindent {\bf Case 1.} There is neither a green out-edge nor a red out-edge on rectangle $A_1$; therefore, rectangle $A_1$ does not have any edges. See Figure \ref{figure-case1-1} for the graph $H$ and Figure \ref{figure-case1-2} for $\|H\|_2^2$. We introduce a partial graph $H'$ which is obtained from $H$ by removing the rectangle $A_1$. See Figure \ref{figure-case1-4} for the graph $H'$ and Figure \ref{figure-case1-3} for $\|H'\|_2^2$.

\begin{figure}[ht]
    \centering
    \begin{subfigure}{0.25\textwidth}
        \centering
        \fitfig{82.25bp}{48.65bp}{%
    \begin{tikzpicture}
      \useasboundingbox (-0.600,-1.237) rectangle (3.605,1.235);
      \grbox{A}{0}{0}{$A_1$}
      \grbox{B}{1.88}{0}{$A_2$}
      \node[dts] at (3.39,0.80) {\dt};
      \node[dts] at (3.39,-0.80) {\dt};
    \end{tikzpicture}%
  }
        \caption{Graph of $H$}
        \label{figure-case1-1}
    \end{subfigure}
    \hspace{-0.5cm}
    \begin{subfigure}{0.25\textwidth}
        \centering
        \fitfig{95.55bp}{125.31bp}{%
    \begin{tikzpicture}
      \useasboundingbox (-1.259,-3.157) rectangle (3.394,3.157);
      \rgbox{Ua}{0}{1.918}{$A_1^{*}$}    \rgbox{Ub}{1.85}{1.918}{$A_2^{*}$}
      \grbox{La}{0}{-1.918}{$A_1$}       \grbox{Lb}{1.85}{-1.918}{$A_2$}
      \draw[gedge] (Uai2) .. controls ($(Uai2)+(-0.45,-0.80)$) and ($(Lai1)+(-0.45,0.80)$) .. (Lai1);
      \draw[gedge] (Lao1) .. controls ($(Lao1)+(0.45,0.80)$)  and ($(Uao2)+(0.45,-0.80)$) .. (Uao2);
      \draw[redge] (Uai1) .. controls ($(Uai1)+(-0.96,-1.70)$) and ($(Lai2)+(-0.96,1.70)$) .. (Lai2);
      \draw[redge] (Lao2) .. controls ($(Lao2)+(0.93,1.70)$)  and ($(Uao1)+(0.93,-1.70)$) .. (Uao1);
      \foreach \y in {2.718,1.118,-1.118,-2.718}{ \node[dts] at (3.22,\y) {\dt}; }
    \end{tikzpicture}%
  }
        \caption{$\|H\|_2^2$}
        \label{figure-case1-2}
    \end{subfigure}
    \hspace{-0.1cm}
    \begin{subfigure}{0.25\textwidth}
        \centering
        \fitfig{43.75bp}{123.56bp}{%
    \begin{tikzpicture}
      \useasboundingbox (-0.579,-3.158) rectangle (1.563,3.158);
      \rgbox{U}{0}{1.918}{$A_2^{*}$}
      \grbox{L}{0}{-1.918}{$A_2$}
      \node[dts] at (1.40,2.718)  {\dt};
      \node[dts] at (1.40,1.118)  {\dt};
      \node[dts] at (1.40,-1.118) {\dt};
      \node[dts] at (1.40,-2.718) {\dt};
    \end{tikzpicture}%
  }
        \caption{$\|H'\|_2^2$}
        \label{figure-case1-3}
    \end{subfigure}
    \hspace{-1cm}
    \begin{subfigure}{0.25\textwidth}
        \centering
        \fitfig{43.05bp}{48.30bp}{%
    \begin{tikzpicture}
      \useasboundingbox (-0.600,-1.237) rectangle (1.576,1.235);
      \grbox{A}{0}{0}{$A_2$}
      \node[dts] at (1.40,0.80) {\dt};
      \node[dts] at (1.40,-0.80) {\dt};
    \end{tikzpicture}%
  }
        \caption{Graph of $H'$}
        \label{figure-case1-4}
    \end{subfigure}
    \caption{Case 1}

\end{figure}

Noting that $A_1^*A_1 = I_N \otimes I_N$, we have
\begin{align} \label{eq-H-norm-case1}
    \|H\|_2^2 = \Tr \left( A_1^* A_1 \right) \|H'\|_2^2 = N^2 \|H'\|_2^2.
\end{align}
Moreover, according to Lemma \ref{Lem-simple-subgraph}, $H'$ is also a simple partial graph of $G_{\sigma_1,\sigma_2}$. The number of rectangles on the graph $H'$ is $n-1$, while the numbers of green edges and red edges are $r$ and $s$, respectively. Thus, by the induction hypothesis, we have $\|H'\|_2^2 = N^{2(n-1)-r-s}$. The proof is concluded by substituting this identity into \eqref{eq-H-norm-case1}.

\bigskip

\noindent {\bf Case 2.} There is no green out-edge, but one red out-edge on rectangle $A_1$. We denote by $A_i$ the successor of $A_1$ with respect to this red out-edge. See Figure \ref{figure-case2-1} for the graph $H$, and Figure \ref{figure-case2-2} for $\|H\|_2^2$. We introduce a partial graph $H'$ that is obtained from $H$ by removing the rectangle $A_1$ and the red out-edge associated with $A_1$. See Figure \ref{figure-case2-4} for the graph $H'$ and Figure \ref{figure-case2-3} for $\|H'\|_2^2$.

\begin{figure}[ht]
    \centering
    \begin{subfigure}{0.25\textwidth}
        \centering
        \fitfig{114.01bp}{47.40bp}{%
    \begin{tikzpicture}
      \useasboundingbox (-0.600,-1.542) rectangle (6.323,1.235);
      \grbox{A}{0}{0}{$A_1$}
      \grbox{B}{1.88}{0}{$A_2$}
      \grbox{C}{4.78}{0}{$A_i$}
      \dnarcs[redge]{Ao2}{Ci2}{0.90}{0.60}
      \node[dts] at (3.39,0.80) {\dt};
      \node[dts] at (3.39,-0.80) {\dt};
      \node[dts] at (6.07,0.80) {\dt};
      \node[dts] at (6.07,-0.80) {\dt};
    \end{tikzpicture}%
  }
        \caption{Graph of $H$}
        \label{figure-case2-1}
    \end{subfigure}
    \hspace{-0.1cm}
    \begin{subfigure}{0.25\textwidth}
        \centering
        \fitfig{109.77bp}{110.33bp}{%
    \begin{tikzpicture}
      \useasboundingbox (-1.166,-3.499) rectangle (5.776,3.459);
      \rgbox{Ua}{0}{1.918}{$A_1^{*}$}  \rgbox{Ub}{1.836}{1.918}{$A_2^{*}$}  \rgbox{Uc}{4.455}{1.918}{$A_i^{*}$}
      \grbox{La}{0}{-1.918}{$A_1$}     \grbox{Lb}{1.836}{-1.918}{$A_2$}     \grbox{Lc}{4.455}{-1.918}{$A_i$}
      \draw[gedge] (Uai2) .. controls ($(Uai2)+(-0.45,-0.80)$) and ($(Lai1)+(-0.45,0.80)$) .. (Lai1);
      \draw[gedge] (Lao1) .. controls ($(Lao1)+(0.45,0.80)$)  and ($(Uao2)+(0.45,-0.80)$) .. (Uao2);
      \draw[redge] (Uai1) .. controls ($(Uai1)+(-0.85,-1.70)$) and ($(Lai2)+(-0.85,1.70)$) .. (Lai2);
      \dnarcs[redge]{Lao2}{Lci2}{0.95}{0.70}
      \uparcs[redge]{Uci1}{Uao1}{0.90}{-0.70}
      \foreach \y in {2.718,1.118,-1.118,-2.718}{
        \node[dts] at (3.22,\y) {\dt};  \node[dts] at (5.60,\y) {\dt}; }
    \end{tikzpicture}%
  }
        \caption{$\|H\|_2^2$}
        \label{figure-case2-2}
    \end{subfigure}
    \begin{subfigure}{0.25\textwidth}
        \centering
        \fitfig{81.00bp}{107.11bp}{%
    \begin{tikzpicture}
      \useasboundingbox (-0.579,-3.158) rectangle (4.221,3.158);
      \rgbox{Ua}{0}{1.918}{$A_2^{*}$}    \rgbox{Ub}{2.93}{1.918}{$A_i^{*}$}
      \grbox{La}{0}{-1.918}{$A_2$}       \grbox{Lb}{2.93}{-1.918}{$A_i$}
      \draw[redge] (Ubi1) .. controls ($(Ubi1)+(-0.885,-1.70)$) and ($(Lbi2)+(-0.885,1.70)$) .. (Lbi2);
      \foreach \y in {2.718,1.118,-1.118,-2.718}{
        \node[dts] at (1.38,\y) {\dt};  \node[dts] at (4.05,\y) {\dt}; }
    \end{tikzpicture}%
  }
        \caption{$\|H'\|_2^2$}
        \label{figure-case2-3}
    \end{subfigure}
    \hspace{-1cm}
    \begin{subfigure}{0.25\textwidth}
        \centering
        \fitfig{70.80bp}{42.90bp}{%
    \begin{tikzpicture}
      \useasboundingbox (-0.600,-1.237) rectangle (3.674,1.235);
      \grbox{A}{0}{0}{$A_2$}
      \grbox{B}{2.48}{0}{$A_i$}
      \node[dts] at (1.39,0.80) {\dt};
      \node[dts] at (1.39,-0.80) {\dt};
      \node[dts] at (3.47,0.80) {\dt};
      \node[dts] at (3.47,-0.80) {\dt};
    \end{tikzpicture}%
  }
        \caption{Graph of $H'$}
        \label{figure-case2-4}
    \end{subfigure}
    \caption{Case 2}

\end{figure}

Since $A_1^*A_1 = I_N \otimes I_N$, we have
\begin{align} \label{eq-H-norm-case2}
    \|H\|_2^2 = \Tr \left( I_N \right) \|H'\|_2^2 = N \|H'\|_2^2.
\end{align}
Moreover, according to Lemma \ref{Lem-simple-subgraph}, $H'$ is also a simple partial graph of $G_{\sigma_1,\sigma_2}$. The number of rectangles in the partial graph $H'$ is $n-1$, and the numbers of green edges and red edges are $r$ and $s-1$, respectively. Thus, by the induction hypothesis, we have $\|H'\|_2^2 = N^{2(n-1)-r-(s-1)}$. The proof is concluded by substituting this identity into \eqref{eq-H-norm-case2}.

\bigskip
\noindent {\bf Case 3.} There is one green out-edge but no red out-edge on rectangle $A_1$. This case is similar to Case 2, and details are omitted.

\bigskip
\noindent {\bf Case 4.} There is one green out-edge and one red out-edge on rectangle $A_1$. See Figure \ref{figure-case4-1} for the graph $H$ and Figure \ref{figure-case4-2} for $\|H\|_2^2$. Let $A_i$ and $A_j$ be the successors of $A_1$ with respect to the red out-edge and the green out-edge, respectively. We introduce a partial graph $H'$ that is obtained from $H$ by removing the rectangle $A_1$ along with the associated red out-edge and green out-edge. See Figure \ref{figure-case4-4} for the graph $H'$ and Figure \ref{figure-case4-3} for $\|H'\|_2^2$.

\begin{figure}[ht]
    \begin{subfigure}{0.48\textwidth}
        \centering
        \fitfig{196.01bp}{70.40bp}{%
    \begin{tikzpicture}
      \useasboundingbox (-0.601,-1.583) rectangle (8.607,1.725);
      \grbox{A}{0}{0}{$A_1$}
      \grbox{B}{1.93}{0}{$A_2$}
      \grbox{C}{4.90}{0}{$A_i$}
      \grbox{D}{7.21}{0}{$A_j$}
      \uparcs[gedge]{Ao1}{Di1}{1.15}{1.80}
      \dnarcs[redge]{Ao2}{Ci2}{0.95}{0.75}
      \foreach \x in {3.47,6.22,8.32}{
        \node[dts] at (\x,0.80) {\dt};  \node[dts] at (\x,-0.80) {\dt}; }
    \end{tikzpicture}%
  }
        \caption{Graph of $H$}
        \label{figure-case4-1}
    \end{subfigure}
    \hfill
    \begin{subfigure}{0.48\textwidth}
        \centering
        \fitfig{193.61bp}{153.61bp}{%
    \begin{tikzpicture}
      \useasboundingbox (-1.166,-3.499) rectangle (7.335,3.425);
      \rgbox{Ua}{0}{1.918}{$A_1^{*}$}  \rgbox{Ub}{1.87}{1.918}{$A_2^{*}$}
      \rgbox{Uc}{3.99}{1.918}{$A_i^{*}$} \rgbox{Ud}{6.07}{1.918}{$A_j^{*}$}
      \grbox{La}{0}{-1.918}{$A_1$}     \grbox{Lb}{1.87}{-1.918}{$A_2$}
      \grbox{Lc}{3.99}{-1.918}{$A_i$}  \grbox{Ld}{6.07}{-1.918}{$A_j$}
      \draw[gedge] (Uai2) .. controls ($(Uai2)+(-0.45,-0.80)$) and ($(Lai1)+(-0.45,0.80)$) .. (Lai1);
      \draw[redge] (Uai1) .. controls ($(Uai1)+(-0.85,-1.70)$) and ($(Lai2)+(-0.85,1.70)$) .. (Lai2);
      \uparcs[gedge]{Lao1}{Ldi1}{1.20}{1.20}
      \dnarcs[gedge]{Udi2}{Uao2}{1.25}{-1.20}
      \dnarcs[redge]{Lao2}{Lci2}{0.95}{0.70}
      \uparcs[redge]{Uci1}{Uao1}{0.85}{-0.70}
      \foreach \y in {2.718,1.118,-1.118,-2.718}{
        \foreach \x in {2.96,5.17,7.17}{ \node[dts] at (\x,\y) {\dt}; } }
    \end{tikzpicture}%
  }
        \caption{$\|H\|_2^2$}
        \label{figure-case4-2}
    \end{subfigure}
    \\
    \begin{subfigure}{0.48\textwidth}
        \centering
        \fitfig{147.21bp}{142.81bp}{%
    \begin{tikzpicture}
      \useasboundingbox (-0.579,-3.158) rectangle (5.998,3.158);
      \rgbox{Ua}{0}{1.918}{$A_2^{*}$}  \rgbox{Ub}{2.57}{1.918}{$A_i^{*}$}
      \rgbox{Uc}{4.73}{1.918}{$A_j^{*}$}
      \grbox{La}{0}{-1.918}{$A_2$}     \grbox{Lb}{2.57}{-1.918}{$A_i$}
      \grbox{Lc}{4.73}{-1.918}{$A_j$}
      \draw[redge] (Ubi1) .. controls ($(Ubi1)+(-0.65,-1.70)$) and ($(Lbi2)+(-0.65,1.70)$) .. (Lbi2);
      \draw[gedge] (Uci2) .. controls ($(Uci2)+(-0.50,-0.80)$) and ($(Lci1)+(-0.50,0.80)$) .. (Lci1);
      \foreach \y in {2.718,1.118,-1.118,-2.718}{
        \foreach \x in {1.17,3.71,5.84}{ \node[dts] at (\x,\y) {\dt}; } }
    \end{tikzpicture}%
  }
        \caption{$\|H'\|_2^2$}
        \label{figure-case4-3}
    \end{subfigure}
    \hfill
    \begin{subfigure}{0.48\textwidth}
        \centering
        \fitfig{131.61bp}{56.40bp}{%
    \begin{tikzpicture}
      \useasboundingbox (-0.600,-1.237) rectangle (5.480,1.235);
      \grbox{A}{0}{0}{$A_2$}
      \grbox{B}{2.11}{0}{$A_i$}
      \grbox{C}{4.13}{0}{$A_j$}
      \foreach \x in {1.02,3.12,5.25}{
        \node[dts] at (\x,0.80) {\dt};  \node[dts] at (\x,-0.80) {\dt}; }
    \end{tikzpicture}%
  }
        \caption{Graph of $H'$}
        \label{figure-case4-4}
    \end{subfigure}
    \caption{Case 4}

\end{figure}

Since $A_1^*A_1 = I_N \otimes I_N$, we have
\begin{align} \label{eq-H-norm-case4}
    \|H\|_2^2 = \|H'\|_2^2.
\end{align}
Moreover, $H'$ is a simple partial graph of the graph $G_{\sigma_1,\sigma_2}$ according to Lemma \ref{Lem-simple-subgraph}. The number of rectangles on the graph $H'$ is $n-1$, and the numbers of green edges and red edges are $r-1$ and $s-1$, respectively. Thus, by the induction hypothesis, we have $\|H'\|_2^2 = N^{2(n-1)-(r-1)-(s-1)} = N^{2n-r-s}$. Together with the identity \eqref{eq-H-norm-case4}, this concludes the proof.
\end{proof}

Now we are ready to prove Theorem \ref{Thm-upper}.

\begin{proof}[Proof of Theorem~\ref{Thm-upper}]
We denote by $G_{\sigma_1,\sigma_2}$ the graph corresponding to \eqref{tr-sigma-tau}. Let $G'_{\sigma_1,\sigma_2}$ be a full simple partial graph of $G_{\sigma_1,\sigma_2}$, and $G''_{\sigma_1,\sigma_2}$ be the complement partial graph of $G'_{\sigma_1,\sigma_2}$. Then, by the Cauchy--Schwarz inequality, we have
\begin{align} \label{eq-inner product-1}
    \left| \left( \Tr_{\sigma_1} \otimes \Tr_{\sigma_2} \right) \left( A_1, \ldots, A_m \right) \right|
    = | G_{\sigma_1,\sigma_2} |
    = \left| \left\langle G'_{\sigma_1,\sigma_2}, G''_{\sigma_1,\sigma_2} \right\rangle \right|
    \le \left\|G'_{\sigma_1,\sigma_2} \right\|_2 \left\| G''_{\sigma_1,\sigma_2} \right\|_2.
\end{align}

We first handle the graph $G''_{\sigma_1,\sigma_2}$. Note that $G'_{\sigma_1,\sigma_2}$ is a full simple partial graph in $G_{\sigma_1,\sigma_2}$, $G''_{\sigma_1,\sigma_2}$ is a graph without rectangles, and the total number of green directed edges and red directed edges is $R_{G'_{\sigma_1,\sigma_2}}$. Thus, we have
\begin{align}\label{eq-norm-G2}
    \left\| G''_{\sigma_1,\sigma_2} \right\|_2^2 = N^{R_{G'_{\sigma_1,\sigma_2}}}.
\end{align}

Next, we turn to the partial graph $G'_{\sigma_1,\sigma_2}$. Note that $G'_{\sigma_1,\sigma_2}$ consists of $m$ rectangles, and the total number of green directed edges and red directed edges is $2m-R_{G'_{\sigma_1,\sigma_2}}$. By Lemma \ref{Lem-inner product}, we have
\begin{align}\label{eq-norm-G1}
    \|G'_{\sigma_1,\sigma_2}\|_2^2
    = N^{2m - \left( 2m-R_{G'_{\sigma_1,\sigma_2}} \right)}
    = N^{R_{G'_{\sigma_1,\sigma_2}}}.
\end{align}

Substituting \eqref{eq-norm-G1} and \eqref{eq-norm-G2} into \eqref{eq-inner product-1}, we obtain
\begin{align*}
    \left| \left( \Tr_{\sigma_1} \otimes \Tr_{\sigma_2} \right) \left( A_1, \ldots, A_m \right) \right|
    \le N^{R_{G'_{\sigma_1,\sigma_2}}}.
\end{align*}

The proof follows by taking the maximum over all unitary matrices $A_1,\ldots,A_m \in M_N (\bC) \otimes M_N(\bC)$, and the minimum over all full simple partial graphs $G'_{\sigma_1,\sigma_2}$ of $G_{\sigma_1,\sigma_2}$.
\end{proof}

\section{Lower bound in the 2 legs case} \label{sec:lower}

In this section, we establish the lower bound for the partial trace \eqref{tr-sigma-tau}.

\subsection{Lower bound for partial trace}

We introduce the following unitary matrix, which plays a key role in the proof.
Let $E_{ij}$ be the $N \times N$ matrix with $1$ in the $(i,j)$ entry and $0$ elsewhere. We let
\begin{align*}
    U = \sum_{i,j=1}^N E_{ij} \otimes E_{ji},
\end{align*}
then $U$ is a unitary matrix in $M_N(\bC) \otimes M_N(\bC)$.
It is actually a self-adjoint involution.

Recall that $M(\sigma_1,\sigma_2)$ is given in Definition \ref{def-cycle-cut}. The following theorem provides a lower bound for the partial trace \eqref{tr-sigma-tau}, and is the main result in this section.

\begin{theorem} \label{Thm-lower}
For $m,N \in \bN$, for any permutations $\sigma_1,\sigma_2 \in \cP([m])$, we have
\begin{align*}
    \max_{A_1,\ldots,A_m \in \{I_N \otimes I_N, U\}} \left| \left( \Tr_{\sigma_1} \otimes \Tr_{\sigma_2} \right) \left( A_1, \ldots, A_m \right) \right|
    = N^{M(\sigma_1,\sigma_2)}.
\end{align*}
\end{theorem}

\begin{proof}
For $\pi\in\cP([2])$, the permutation tensor $E_\pi$ of \eqref{eq-permutation-tensor} is $I_N\otimes I_N$ if $\pi=(1)(2)$, and $U$ if $\pi=(12)$. Thus, in the graph $G_{\sigma_1,\sigma_2}$, the two blue directed edges in each rectangle must either connect vertices of the same color or connect vertices of different colors. They correspond to $I_N\otimes I_N$ and $U$, respectively, as shown in Figure~\ref{figure-blue edges-U-I}.
\begin{figure}[ht]
    \begin{subfigure}{0.48\textwidth}
        \centering
        \fitfig{28.80bp}{56.00bp}{%
    \begin{tikzpicture}
      \useasboundingbox (-0.600,-1.237) rectangle (0.578,1.235);
      \grbox{A}{0}{0}{$A_1$}
      \draw[bedge] (Ai1) -- (Ao1);
      \draw[bedge] (Ai2) -- (Ao2);
    \end{tikzpicture}%
  }
        \caption{$A_1 = I_N \otimes I_N$}
    \end{subfigure}
    \hfill
    \begin{subfigure}{0.48\textwidth}
        \centering
        \fitfig{25.60bp}{57.20bp}{%
    \begin{tikzpicture}
      \useasboundingbox (-0.600,-1.237) rectangle (0.578,1.235);
      \grbox{A}{0}{0}{$A_1$}
      \draw[bedge] (Ai1) -- (Ao2);
      \draw[bedge] (Ai2) -- (Ao1);
    \end{tikzpicture}%
  }
        \caption{$A_1 = U$}
    \end{subfigure}
    \caption{Correspondence of blue directed edges and matrix}
    \label{figure-blue edges-U-I}
\end{figure}

For a fixed blue-edge configuration $\boldsymbol\pi=(\pi_1,\ldots,\pi_m)$ for the graph $G_{\sigma_1,\sigma_2}$, equation~\eqref{eq-coordinate-cycles} gives
\begin{equation}
(\Tr_{\sigma_1}\otimes\Tr_{\sigma_2})(E_{\pi_1},\ldots,E_{\pi_m})
=N^{M(\sigma_1,\sigma_2;\boldsymbol\pi)}.
\end{equation}
Each directed cycle in the graph $G_{\sigma_1,\sigma_2}$ contributes one independent index ranging over $[N]$. Every tuple from $\{I_N\otimes I_N,U\}^m$ is realized by such a blue-edge configuration, so its value is at most $N^{M(\sigma_1,\sigma_2)}$. Choosing a blue-edge configuration with the maximum number of cycles attains this value and proves the theorem.
\end{proof}

\section{Proof of Theorem \ref{Thm-main} and Corollary \ref{Coro-main-2 leg}} \label{sec:2legs}

In this section, we first develop the proof of Theorem~\ref{Thm-main} with the help of Theorems~\ref{Thm-upper} and~\ref{Thm-lower}.


\begin{proof}[Proof of Theorem~\ref{Thm-main}]
Since $I_N\otimes I_N$ and $U$ are unitary, Theorem~\ref{Thm-lower} gives the lower bound over all unitary tuples, and Theorem~\ref{Thm-upper} gives the upper bound. Hence
\begin{equation}
N^{M(\sigma_1,\sigma_2)}
\le \max_{A_i\text{ unitary}}|(\Tr_{\sigma_1}\otimes\Tr_{\sigma_2})(A_1,\ldots,A_m)|
\le N^{C(\sigma_1,\sigma_2)}.
\end{equation}
\end{proof}

\end{document}
